\documentclass{article}
\usepackage[T1]{fontenc}
\usepackage{lmodern}
\usepackage{graphicx}
\usepackage{amsmath,amssymb}
\usepackage[a4paper,margin=2cm]{geometry}
\usepackage[absolute,overlay]{textpos}
\setlength{\TPHorizModule}{1mm}
\setlength{\TPVertModule}{1mm}
\usepackage[indonesian]{babel}
\usepackage{hyperref}
\setlength{\emergencystretch}{2em}
% unit: Halaman 48

\begin{document}

% === Halaman 48 (python/native) ===

• Kelemahan Double DES: serangan meet-in-the-middle attack:

• Dari pengamatan, 


C = EK2(EK1(P))

   maka 


X = EK1(P) = DK2(C)


• Misalkan kriptanalis memiliki potongan C dan P yang berkorepsonden (known-

plaintext attack).

• Enkripsi P untuk semua kemungkinan nilai K1 (yaitu sebanyak 256 kemungkinan 

kunci). Hasilnya adalah semua nilai X  

• Simpan semua nilai X ini di dalam tabel 

48

E

E

X

P

C

K1

K2

Enkripsi

D

D

X

C

P

K2

K1

Dekripsi

\end{document}
